\section{Relation to spectral and coefficient methods}
\label{app:spectral_context}
Finite Fourier
representations of quantum features are established tools
\cite{schuldFourier}; the factorization $C^TMC$ is also the familiar
pullback of a Mahalanobis metric \cite{weinberger}.
Our use is an exact, readout-specific parameter-identifiability
contract and its application to the rank gaps below, not a claim
that either algebraic construction is new.

Fourier coefficient correlations already characterize ansatz-dependent
constraints \cite{fingerprints}. Spectral tensor-network work represents
quantum kernels by Fourier coefficient tensors and uses an empirical
Fourier Gram matrix for data-dependent alignment \cite{spectralTN}.
That data matrix is not the parameter-dependent $C(\theta)^TMC(\theta)$
here, and tensor bond rank is not the differential rank in
Eq.~\eqref{eq:gramrank}. Fourier-geometric gate-placement criteria
also exploit adjoint frequency planes \cite{fourierDesign}.
These overlaps motivate careful separation of representation,
coefficient reachability, and exact parameter identifiability.
Campbell, Del Debbio, and Wallden's circuit harmonic matrix also
eliminates data dependence from coefficient-space analysis
\cite{CHM}. It represents joint input--parameter harmonics and
constructs pointwise coefficient Jacobians and QNTKs, as well as
parameter-averaged statistics. Our distinction is therefore not
``pointwise versus averaged'': $G=C(\theta)^TMC(\theta)$ describes
fixed-metric feature distances, whereas their harmonic QNTK is a
Gram matrix of coefficient derivatives. We use this distinction to
certify distance-equivalence fibers and their differential dimension.
In the four-qubit Hubregtsen example of Sec.~\ref{sec:gram}, the
coefficient Jacobian has rank four, whereas $D_\theta\operatorname{vech}G$
has rank two.
